\documentclass[aspectratio=169,10pt]{beamer}
\makeatletter\def\input@path{{../estilo/}{estilo/}}\makeatother
\usepackage{devmobile}
\modulo{05}{Estado e Navegação}{Estado em Compose, recomposição, state hoisting, efeitos colaterais, ciclo de vida e Navigation 3}

\begin{document}
\capa

\begin{frame}{Objetivos do módulo}
  \begin{conceito}[Ao final deste módulo você será capaz de]
    \begin{itemize}
      \item Criar estado com \cod{remember} e \cod{mutableStateOf} e entender a recomposição.
      \item Preservar estado em rotações com \cod{rememberSaveable}.
      \item Aplicar \emph{state hoisting} para criar componentes reutilizáveis e testáveis.
      \item Usar efeitos colaterais com segurança (\cod{LaunchedEffect}, \cod{DisposableEffect}\ldots).
      \item Explicar o ciclo de vida da Activity e coletar fluxos respeitando-o.
      \item Implementar navegação entre telas com \textbf{Navigation 3}.
    \end{itemize}
  \end{conceito}
\end{frame}

\sumario

% =============================================================================
\section{Estado e recomposição}

\begin{frame}[fragile]{O que é estado}
  \begin{columns}[T]
    \begin{column}{0.5\textwidth}
      \small
      \termo{Estado} é qualquer valor que pode mudar com o tempo e afeta o que aparece na tela: o texto digitado, a lista carregada, se um diálogo está aberto, a última leitura do sensor.
\begin{codigo}[title=Não funciona!]{Kotlin}
@Composable
fun Contador() {
  var cliques = 0          // variável comum
  Button(onClick = { cliques++ }) {
    Text("Cliques: $cliques")  // nunca muda
  }
}
\end{codigo}
      \begin{erro}
        Mudar uma variável comum não avisa o Compose, e a cada recomposição ela volta a 0.
      \end{erro}
    \end{column}
    \begin{column}{0.48\textwidth}
\begin{codigo}[title=Correto]{Kotlin}
@Composable
fun Contador() {
  var cliques by remember { mutableStateOf(0) }
  Button(onClick = { cliques++ }) {
    Text("Cliques: $cliques")
  }
}
\end{codigo}
      \begin{conceito}[Duas peças]
        \begin{itemize}
          \item \cod{mutableStateOf}: valor \textbf{observável}; ler registra dependência, escrever agenda recomposição.
          \item \cod{remember}: guarda o objeto entre recomposições.
        \end{itemize}
      \end{conceito}
    \end{column}
  \end{columns}
\end{frame}

\begin{frame}{Recomposição}
  \begin{columns}[c]
    \begin{column}{0.52\textwidth}
      \centering
      \begin{tikzpicture}[node distance=4mm]
        \node[caixaescura,minimum width=4.6cm] (c) {\textbf{Composição inicial}\\ executa as funções e cria a árvore};
        \node[caixaambar,minimum width=4.6cm,below=of c] (m) {\textbf{Estado muda}\\ \cod{cliques++}};
        \node[caixa,minimum width=4.6cm,below=of m] (r) {\textbf{Recomposição}\\ reexecuta \emph{só} quem leu o estado};
        \node[caixa,minimum width=4.6cm,below=of r,fill=cinzaclaro] (d) {layout e desenho atualizados};
        \draw[seta] (c) -- (m); \draw[seta] (m) -- (r); \draw[seta] (r) -- (d);
        \draw[setaambar,rounded corners] (d.east) -- ++(5mm,0) |- (m.east);
      \end{tikzpicture}
    \end{column}
    \begin{column}{0.46\textwidth}
      \begin{conceito}[Propriedades]
        \begin{itemize}
          \item É \textbf{inteligente}: funções cujos parâmetros não mudaram podem ser puladas.
          \item Pode acontecer \textbf{a cada quadro} (ex.: animações).
          \item Pode ser \textbf{cancelada} e refeita.
        \end{itemize}
      \end{conceito}
      \begin{dica}
        O \emph{Layout Inspector} mostra quantas vezes cada composable recompôs. Números muito altos indicam problema de desempenho (Módulo 11).
      \end{dica}
    \end{column}
  \end{columns}
\end{frame}

\begin{frame}[fragile]{\cod{remember} \emph{versus} \cod{rememberSaveable}}
  \begin{columns}[T]
    \begin{column}{0.52\textwidth}
\begin{codigo}[title=Sobrevivendo à rotação]{Kotlin}
// perdido ao girar a tela
var filtro by remember { mutableStateOf("") }

// salvo no Bundle: sobrevive à rotação e à
// morte do processo em segundo plano
var busca by rememberSaveable { mutableStateOf("") }

// estado derivado: só recalcula quando muda
val visiveis by remember(itens) {
  derivedStateOf { itens.filter { busca in it.descricao } }
}

// listas observáveis
val selecionados = remember { mutableStateListOf<String>() }
\end{codigo}
    \end{column}
    \begin{column}{0.46\textwidth}
      \footnotesize
      \begin{tabularx}{\linewidth}{@{}l c c@{}}
        \toprule
        & \cabecalho{recomposição} & \cabecalho{rotação} \\
        \midrule
        variável comum & \textcolor{vermelhoerro}{perde} & \textcolor{vermelhoerro}{perde} \\
        \cod{remember} & \textcolor{verdeok}{mantém} & \textcolor{vermelhoerro}{perde} \\
        \cod{rememberSaveable} & \textcolor{verdeok}{mantém} & \textcolor{verdeok}{mantém} \\
        ViewModel & \textcolor{verdeok}{mantém} & \textcolor{verdeok}{mantém} \\
        \bottomrule
      \end{tabularx}
      \begin{atencao}
        \cod{rememberSaveable} só aceita tipos que cabem num \cod{Bundle} (primitivos, \cod{String}, \cod{Parcelable}). Para objetos maiores, use um ViewModel (Módulo 06).
      \end{atencao}
    \end{column}
  \end{columns}
\end{frame}

% =============================================================================
\section{State hoisting}

\begin{frame}[fragile]{Elevando o estado (\emph{state hoisting})}
  \begin{columns}[T]
    \begin{column}{0.52\textwidth}
\begin{codigo}[title=Componente sem estado (stateless)]{Kotlin}
@Composable
fun CampoLimite(
  valor: String,                   // estado desce
  onValorChange: (String) -> Unit, // evento sobe
  modifier: Modifier = Modifier
) {
  OutlinedTextField(
    value = valor,
    onValueChange = onValorChange,
    label = { Text("Limite (°C)") },
    modifier = modifier
  )
}

// Quem usa decide onde o estado mora
@Composable
fun TelaConfig() {
  var limite by rememberSaveable { mutableStateOf("30") }
  CampoLimite(limite, onValorChange = { limite = it })
}
\end{codigo}
    \end{column}
    \begin{column}{0.46\textwidth}
      \centering
      \begin{tikzpicture}[node distance=12mm]
        \node[caixaescura,minimum width=3.6cm] (p) {\cod{TelaConfig}\\ \tiny dona do estado};
        \node[caixa,minimum width=3.6cm,below=of p] (f) {\cod{CampoLimite}\\ \tiny stateless};
        \draw[seta] ([xshift=-8mm]p.south) -- node[rotulo,left]{estado} ([xshift=-8mm]f.north);
        \draw[setaambar] ([xshift=8mm]f.north) -- node[rotulo,right]{eventos} ([xshift=8mm]p.south);
      \end{tikzpicture}
      \begin{conceito}[Fluxo unidirecional de dados (UDF)]
        O estado desce como parâmetro e os eventos sobem como lambdas. Há uma única \textbf{fonte da verdade}.
      \end{conceito}
      \begin{dica}[Vantagens]
        Componente reutilizável, fácil de testar e de visualizar em \cod{@Preview}.
      \end{dica}
    \end{column}
  \end{columns}
\end{frame}

\begin{frame}{Onde o estado deve morar?}
  \begin{center}
  \begin{tikzpicture}[node distance=5mm and 8mm,
      decisao/.style={caixaambar,diamond,aspect=2.6,inner sep=1pt,text width=2.5cm,font=\tiny}]
    \node[decisao] (q1) {É só aparência da UI? (expandido, rolagem, foco)};
    \node[caixaescura,right=12mm of q1,text width=3.2cm] (r1) {No próprio composable,\\ com \cod{remember(Saveable)}};
    \node[decisao,below=6mm of q1] (q2) {Vários composables usam?};
    \node[caixa,right=12mm of q2,text width=3.2cm] (r2) {No ancestral comum mais baixo (hoisting)};
    \node[decisao,below=6mm of q2] (q3) {Envolve regras de negócio ou dados?};
    \node[caixaambarforte,right=12mm of q3,text width=3.2cm] (r3) {No \textbf{ViewModel} (\emph{state holder} da tela)};
    \draw[seta] (q1) -- node[rotulo,above]{sim} (r1);
    \draw[seta] (q1) -- node[rotulo,left]{não} (q2);
    \draw[seta] (q2) -- node[rotulo,above]{sim} (r2);
    \draw[seta] (q2) -- node[rotulo,left]{não} (q3);
    \draw[seta] (q3) -- node[rotulo,above]{sim} (r3);
  \end{tikzpicture}
  \end{center}
  \begin{dica}
    Regra prática: eleve o estado só até onde for necessário, nem mais nem menos.
  \end{dica}
\end{frame}

% =============================================================================
\section{Efeitos colaterais}

\begin{frame}[fragile]{APIs de efeitos}
  \begin{columns}[T]
    \begin{column}{0.5\textwidth}
\begin{codigo}[title=LaunchedEffect e DisposableEffect]{Kotlin}
@Composable
fun Monitor(bancadaId: Int, vm: PainelViewModel) {
  // roda quando entra na tela e quando a chave muda
  LaunchedEffect(bancadaId) {
    vm.carregar(bancadaId)   // coroutine
  }

  val lifecycle = LocalLifecycleOwner.current.lifecycle
  DisposableEffect(lifecycle) {
    val obs = LifecycleEventObserver { _, e ->
      if (e == Lifecycle.Event.ON_RESUME) vm.reconectar()
    }
    lifecycle.addObserver(obs)
    onDispose { lifecycle.removeObserver(obs) }
  }
}
\end{codigo}
    \end{column}
    \begin{column}{0.48\textwidth}
      \scriptsize
      \begin{tabularx}{\linewidth}{@{}l L@{}}
        \toprule
        \cabecalho{API} & \cabecalho{Quando usar} \\
        \midrule
        \cod{LaunchedEffect(chave)} & rodar \cod{suspend} ao entrar ou quando a chave muda \\
        \cod{rememberCoroutineScope} & iniciar coroutine a partir de um clique \\
        \cod{DisposableEffect} & registrar e remover \emph{listeners} \\
        \cod{SideEffect} & sincronizar com código não-Compose a cada recomposição \\
        \cod{rememberUpdatedState} & usar o valor mais recente dentro de um efeito longo \\
        \cod{produceState} & converter fonte externa em estado \\
        \cod{derivedStateOf} & estado calculado que muda menos que a origem \\
        \bottomrule
      \end{tabularx}
      \begin{atencao}
        \cod{LaunchedEffect(Unit)} roda uma única vez enquanto o composable estiver na tela.
      \end{atencao}
    \end{column}
  \end{columns}
\end{frame}

% =============================================================================
\section{Ciclo de vida}

\begin{frame}{Ciclo de vida da Activity}
  \begin{columns}[c]
    \begin{column}{0.56\textwidth}
      \centering
      \begin{tikzpicture}[node distance=2.2mm,cb/.style={caixa,minimum width=2.3cm,minimum height=5.5mm,font=\tiny}]
        \node[caixaambarforte,minimum width=2.3cm] (l) {Lançada};
        \node[cb,below=of l] (c) {\cod{onCreate()}};
        \node[cb,below=of c] (s) {\cod{onStart()}};
        \node[cb,below=of s] (r) {\cod{onResume()}};
        \node[caixaescura,minimum width=2.3cm,below=of r] (run) {Em execução};
        \node[cb,below=of run] (p) {\cod{onPause()}};
        \node[cb,below=of p] (st) {\cod{onStop()}};
        \node[cb,below=of st] (d) {\cod{onDestroy()}};
        \foreach \a/\b in {l/c,c/s,s/r,r/run,run/p,p/st,st/d}{\draw[seta] (\a) -- (\b);}
        \node[cb,right=7mm of s] (rs) {\cod{onRestart()}};
        \draw[setaambar] (st.east) -| (rs.south);
        \draw[setaambar] (rs.west) -- (s.east);
        \draw[setaambar] (p.west) -- ++(-6mm,0) |- (r.west);
        \node[rotulo,anchor=east] at ([xshift=-7mm]$(p.west)!0.5!(r.west)$) {volta ao\\ primeiro plano};
        \node[rotulo,anchor=north] at ([yshift=-1mm]rs.south east) {volta do\\ segundo plano};
      \end{tikzpicture}
    \end{column}
    \begin{column}{0.42\textwidth}
      \scriptsize
      \begin{tabularx}{\linewidth}{@{}l L@{}}
        \toprule
        \cabecalho{Estado} & \cabecalho{Significado} \\
        \midrule
        \textbf{Created} & criada, ainda invisível \\
        \textbf{Started} & visível \\
        \textbf{Resumed} & em primeiro plano, recebe toques \\
        \textbf{Paused} & parcialmente encoberta \\
        \textbf{Stopped} & invisível (segundo plano) \\
        \bottomrule
      \end{tabularx}
      \begin{atencao}[Mudanças de configuração]
        Girar a tela, trocar idioma ou tema \textbf{destrói e recria} a Activity. Por isso o estado precisa estar em \cod{rememberSaveable} ou no ViewModel.
      \end{atencao}
      \begin{atencao}[Morte do processo]
        Em segundo plano, o sistema pode encerrar o app para liberar memória. Ao voltar, só o estado salvo sobrevive.
      \end{atencao}
    \end{column}
  \end{columns}
\end{frame}

\begin{frame}[fragile]{Coletando fluxos com segurança}
  \begin{columns}[T]
    \begin{column}{0.54\textwidth}
\begin{codigo}[title=TelaPainel.kt]{Kotlin}
@Composable
fun TelaPainel(vm: PainelViewModel = viewModel()) {
  // para de coletar quando o app vai para o
  // segundo plano (abaixo de STARTED)
  val estado by vm.estado.collectAsStateWithLifecycle()

  when (val e = estado) {
    PainelUiState.Carregando -> CircularProgressIndicator()
    is PainelUiState.Erro -> Text(e.mensagem)
    is PainelUiState.Pronto -> Painel(
      leituras = e.leituras,
      onAtualizar = vm::atualizar
    )
  }
}
\end{codigo}
    \end{column}
    \begin{column}{0.44\textwidth}
      \begin{conceito}[\cod{collectAsStateWithLifecycle}]
        Converte \cod{StateFlow} em \cod{State} do Compose e \textbf{pausa a coleta} quando a tela não está visível, economizando bateria, rede e sensores.
      \end{conceito}
      \begin{dica}[Dependência]
        Vem do artefato\\ \cod{androidx.lifecycle:}\\\cod{lifecycle-runtime-compose}.
      \end{dica}
      \begin{erro}
        \cod{collectAsState()} continua coletando em segundo plano. No Android, prefira a versão \emph{WithLifecycle}.
      \end{erro}
    \end{column}
  \end{columns}
\end{frame}

% =============================================================================
\section{Navegação com Navigation 3}

\begin{frame}{Conceitos de navegação}
  \begin{columns}[c]
    \begin{column}{0.5\textwidth}
      \small
      \begin{itemize}
        \item Apps Compose costumam ter \textbf{uma Activity} e várias telas (\emph{destinos}).
        \item A \termo{pilha de navegação} (\emph{back stack}) guarda o histórico; o botão/gesto voltar remove o topo.
        \item \textbf{Navigation 3} (Jetpack, estável desde nov./2025) representa a pilha como uma \textbf{lista observável de chaves} controlada pelo próprio app.
      \end{itemize}
      \begin{conceito}[Peças]
        \cod{NavKey} (chave de cada tela) $\cdot$ \cod{rememberNavBackStack} (a pilha) $\cdot$ \cod{NavDisplay} (exibe o topo) $\cdot$ \cod{entryProvider} (chave $\rightarrow$ tela).
      \end{conceito}
    \end{column}
    \begin{column}{0.48\textwidth}
      \centering
      \begin{tikzpicture}[node distance=1mm]
        \node[caixa,minimum width=3.2cm] (a) {\cod{Lista}};
        \node[caixa,minimum width=3.2cm,above=of a] (b) {\cod{Detalhe(id = 42)}};
        \node[caixaambarforte,minimum width=3.2cm,above=of b] (c) {\cod{Reserva(id = 42)}};
        \node[rotulo,right=2mm of c] {topo: visível};
        \draw[setaambar] ([xshift=-4mm]c.north west) -- ++(0,4mm) node[above,rotulo]{\cod{add()}};
        \draw[seta] ([xshift=4mm]c.north east) ++(0,4mm) -- ([xshift=4mm]c.north east) ;
        \node[rotulo] at ([xshift=12mm,yshift=6mm]c.north east) {\cod{removeLastOrNull()}};
        \node[pontilhado,fit=(a)(b)(c)] {};
      \end{tikzpicture}
      \begin{dica}
        Muitos projetos existentes usam \textbf{Navigation Compose} (Nav2, com \cod{NavHost} e \cod{NavController}). Os conceitos são os mesmos; ao final mostramos a equivalência.
      \end{dica}
    \end{column}
  \end{columns}
\end{frame}

\begin{frame}[fragile]{Navigation 3: chaves e pilha}
  \begin{columns}[T]
    \begin{column}{0.5\textwidth}
\begin{codigo}[title=app/build.gradle.kts]{Kotlin}
plugins {
  alias(libs.plugins.kotlin.serialization)
}
dependencies {
  implementation(libs.androidx.navigation3.runtime)
  implementation(libs.androidx.navigation3.ui)
  implementation(
    libs.androidx.lifecycle.viewmodel.navigation3)
  implementation(libs.kotlinx.serialization.core)
}
\end{codigo}
\begin{codigo}[title=Rotas.kt]{Kotlin}
@Serializable data object Lista : NavKey
@Serializable data class Detalhe(val id: Int) : NavKey
@Serializable data class Reserva(val id: Int) : NavKey
\end{codigo}
    \end{column}
    \begin{column}{0.48\textwidth}
      \begin{conceito}[Chaves tipadas]
        Cada tela é uma classe \cod{@Serializable} que implementa \cod{NavKey}. Os argumentos são propriedades: sem strings de rota e sem erros de digitação.
      \end{conceito}
      \begin{conceito}[Pilha persistente]
        \cod{rememberNavBackStack(Lista)} cria a pilha com a tela inicial e a salva automaticamente em rotações e na morte do processo, por isso as chaves são serializáveis.
      \end{conceito}
      \begin{dica}
        Passe só identificadores (\cod{id}) nas chaves; a tela de destino busca os dados.
      \end{dica}
    \end{column}
  \end{columns}
\end{frame}

\begin{frame}[fragile]{Navigation 3: \cod{NavDisplay}}
  \begin{columns}[T]
    \begin{column}{0.6\textwidth}
\begin{codigo}[title=NavegacaoApp.kt]{Kotlin}
@Composable
fun NavegacaoApp() {
  val pilha = rememberNavBackStack(Lista)
  NavDisplay(
    backStack = pilha,
    onBack = { pilha.removeLastOrNull() },
    entryDecorators = listOf(
      rememberSaveableStateHolderNavEntryDecorator(),
      rememberViewModelStoreNavEntryDecorator()), // VM por tela
    entryProvider = entryProvider {
      entry<Lista> {
        TelaLista(onAbrir = { id -> pilha.add(Detalhe(id)) })
      }
      entry<Detalhe> { chave ->
        TelaDetalhe(chave.id,
          onReservar = { pilha.add(Reserva(chave.id)) })
      }
      entry<Reserva> { chave -> TelaReserva(chave.id) }
    }
  )
}
\end{codigo}
    \end{column}
    \begin{column}{0.38\textwidth}
      \small
      \begin{itemize}
        \item \textbf{Navegar} = \cod{pilha.add(chave)}.
        \item \textbf{Voltar} = \cod{pilha.removeLastOrNull()}.
        \item O gesto de voltar do sistema chama \cod{onBack}.
        \item Os decoradores salvam o estado de cada tela e criam um ViewModel por entrada da pilha.
      \end{itemize}
      \begin{atencao}[Telas não conhecem a pilha]
        \cod{TelaLista} recebe \cod{onAbrir} como lambda, não a pilha. Assim ela continua reutilizável e testável.
      \end{atencao}
    \end{column}
  \end{columns}
\end{frame}

\begin{frame}[fragile]{Navegação principal com barra inferior}
  \begin{columns}[T]
    \begin{column}{0.56\textwidth}
\begin{codigo}[title=Abas principais]{Kotlin}
enum class Aba(val chave: NavKey, val rotulo: String,
               val icone: ImageVector) {
  PAINEL(Painel, "Painel", Icons.Default.Dashboard),
  EQUIP(Lista, "Equipamentos", Icons.Default.Memory),
  CONFIG(Config, "Ajustes", Icons.Default.Settings)
}

NavigationBar {
  Aba.entries.forEach { aba ->
    NavigationBarItem(
      selected = pilha.first() == aba.chave,
      onClick = {           // troca a raiz da pilha
        pilha.clear(); pilha.add(aba.chave)
      },
      icon = { Icon(aba.icone, null) },
      label = { Text(aba.rotulo) }
    )
  }
}
\end{codigo}
    \end{column}
    \begin{column}{0.42\textwidth}
      \begin{dica}[Uma pilha por aba]
        Em apps maiores, mantenha uma pilha para cada aba e exiba a da aba ativa. Assim cada aba lembra até onde o usuário navegou. Os \emph{recipes} oficiais do Navigation 3 trazem esse padrão pronto.
      \end{dica}
      \begin{conceito}[Nav2 $\leftrightarrow$ Nav3]
        \scriptsize
        \begin{tabular}{@{}l l@{}}
          \cod{NavHost} & \cod{NavDisplay} \\
          \cod{navigate(r)} & \cod{pilha.add(r)} \\
          \cod{popBackStack()} & \cod{removeLastOrNull()} \\
          \cod{composable<R>\{\}} & \cod{entry<R>\{\}} \\
          \cod{toRoute<R>()} & a própria chave \\
        \end{tabular}
      \end{conceito}
    \end{column}
  \end{columns}
\end{frame}

% =============================================================================
\section{Fechamento}

\begin{frame}{Resumo}
  \begin{resumo}
    \begin{itemize}
      \item \cod{mutableStateOf} + \cod{remember} criam estado observável; mudar o estado dispara a recomposição.
      \item \cod{rememberSaveable} sobrevive à rotação; o ViewModel guarda o estado da tela.
      \item \emph{State hoisting}: estado desce, eventos sobem, uma fonte da verdade.
      \item Efeitos colaterais só dentro de \cod{LaunchedEffect}, \cod{DisposableEffect} e afins.
      \item Rotação recria a Activity; colete fluxos com \cod{collectAsStateWithLifecycle}.
      \item Navigation 3: chaves \cod{NavKey}, pilha observável e \cod{NavDisplay}.
    \end{itemize}
  \end{resumo}
\end{frame}

\begin{frame}{Exercícios}
  \begin{exercicio}[Exercício 1 --- Rotação]
    Crie um formulário com três campos, um usando \cod{remember} e dois com \cod{rememberSaveable}. Gire o emulador e explique o que aconteceu.
  \end{exercicio}
  \begin{exercicio}[Exercício 2 --- Hoisting]
    Transforme um \cod{SeletorDeFaixa} com estado interno (mínimo e máximo) em um componente \emph{stateless} com \cod{RangeSlider}. Escreva dois \cod{@Preview}s com valores diferentes.
  \end{exercicio}
  \begin{exercicio}[Exercício 3 --- Ciclo de vida]
    Registre um \cod{Log} em cada callback do ciclo de vida da \cod{MainActivity}. Gire a tela, vá para a tela inicial e volte. Desenhe a sequência observada.
  \end{exercicio}
  \begin{projeto}[LabMonitor --- Etapa 5]
    Implemente a navegação: \textbf{Painel}, \textbf{Equipamentos} (lista $\rightarrow$ detalhe $\rightarrow$ reserva) e \textbf{Ajustes} com \cod{NavigationBar}. O formulário de reserva deve sobreviver à rotação.
  \end{projeto}
\end{frame}

\begin{frame}{Para aprofundar}
  \small
  \begin{itemize}
    \item Estado e Jetpack Compose: \url{https://developer.android.com/develop/ui/compose/state}
    \item Onde elevar o estado: \url{https://developer.android.com/develop/ui/compose/state-hoisting}
    \item Efeitos colaterais: \url{https://developer.android.com/develop/ui/compose/side-effects}
    \item Ciclo de vida da Activity: \url{https://developer.android.com/guide/components/activities/activity-lifecycle}
    \item Navigation 3: \url{https://developer.android.com/guide/navigation/navigation-3}
    \item Recipes do Navigation 3: \url{https://github.com/android/nav3-recipes}
  \end{itemize}
\end{frame}

\encerramento{Módulo 06 --- Arquitetura de apps: MVVM, camadas, repositório, casos de uso e injeção de dependência com Hilt.}

\end{document}
